\chapter{Foundational principles}
\label{chap:principi-fondanti}

\section{Zero-trust}
\label{sec:sota-zero-trust}

The term ``zero trust'' was coined in 2010 by Forrester analyst John Kindervag. The underlying idea
was not to grant trust on the basis of network location, but to verify every access
independently of it.

NIST Special Publication 800-207 \cite{nist-sp-800-207} today formalizes this principle as the assumption that no network,
internal or external, should be considered implicitly trustworthy, and that every access request
should be evaluated independently of its origin.

RAM-USB extends the zero-trust principle by distrusting even the identity of a component authenticated via mTLS.
This is because every component assumes that the others are under the control of an attacker.


\section{Defense-in-depth}
\label{sec:sota-defence-in-depth}

NIST Special Publication 800-53 Revision 5 defines defense-in-depth as a strategy that integrates
people, technology, and operational capabilities to establish redundant barriers across multiple layers of
an organization, so that the failure or circumvention of a single barrier does not compromise the system
as a whole \cite{nist-sp-800-53r5}.
